%% $Id: pst-vehicle.tex 376 2021-12-29 12:28:42Z herbert $
%%
%% This is file `pst-vehicle.tex',
%%
%% IMPORTANT NOTICE:
%%
%% Package `pst-vehicle.tex'
%%
%% Thomas Söll
%% with the collaboration of
%% Juergen Gilg
%% Manuel Luque
%% Herbert Voß (bugfixes)
%%
%% This program can redistributed and/or modified under %%
%% the terms of the LaTeX Project Public License        %%
%% Distributed from CTAN archives in directory          %%
%% macros/latex/base/lppl.txt; either version 1.3c of   %%
%% the License, or (at your option) any later version.  %%
%%
%% DESCRIPTION:
%%   `pst-vehicle' is a PSTricks package
%%
%%
\csname PSTvehicleLoaded\endcsname
\let\PSTvehicleLoaded\endinput
% Requires PSTricks, pst-xkey, pst-node packages
\ifx\PSTricksLoaded\endinput\else\input pstricks.tex\fi
\ifx\PSTXKeyLoaded\endinput\else\input pst-xkey.tex\fi
\ifx\PSTplotLoaded\endinput\else\input pst-plot.tex\fi
\ifx\PSTnodeLoaded\endinput\else\input pst-node.tex\fi
\def\fileversion{1.3}
\def\filedate{2021/12/29}
\message{`PST' v\fileversion, \filedate}

\edef\PstAtCode{\the\catcode`\@} \catcode`\@=11\relax

\definecolor{greenSlp}{rgb}{0,0.4,0.0}
\definecolor{redSlp}{rgb}{0.7,0.134,0.134}

\pst@addfams{pst-vehicle}
\define@key[psset]{pst-vehicle}{backwheel}[\wheelA]{\def\pst@backwheel{#1 }} %
\define@key[psset]{pst-vehicle}{frontwheel}[\wheelA]{\def\pst@frontwheel{#1 }} %
\define@key[psset]{pst-vehicle}{ownvehicle}{\def\pst@ownvehicle{#1 }} %
\define@key[psset]{pst-vehicle}{vehicle}[\Bike]{\def\pst@vehicle{#1}} %
\define@key[psset]{pst-vehicle}{GravNode}[]{\def\pst@GravNode{#1 }} %
\define@key[psset]{pst-vehicle}{d}[5.8]{\def\pst@d{#1 }} %
\define@key[psset]{pst-vehicle}{rF}[1.6]{\def\pst@rF{#1 }} %
\define@key[psset]{pst-vehicle}{rB}[1.6]{\def\pst@rB{#1 }} %
\define@key[psset]{pst-vehicle}{gang}[1]{\def\pst@gang{#1 }} %
\define@key[psset]{pst-vehicle}{epsilon}[1e-6]{\def\pst@epsilon{#1 }} %
\define@key[psset]{pst-vehicle}{startPos}[0]{\def\pst@startPos{#1 }} %
\define@boolkey[psset]{pst-vehicle}[Pst@]{showSlope}[true]{} %
\define@boolkey[psset]{pst-vehicle}[Pst@]{MonoAxis}[false]{} %
%\define@boolkey[psset]{pst-vehicle}[Pst@]{EqValveStartPos}[true]{} %

%% $Id: pst-vehicle.tex 376 2021-12-29 12:28:42Z herbert $
%%
%% This is file `pst-vehicle.tex',
%%
%% IMPORTANT NOTICE:
%%
%% Package `pst-vehicle.tex'
%%
%% Thomas Söll
%% with the collaboration of
%% Juergen Gilg
%% Manuel Luque
%% Herbert Voß (bugfixes)
%%
%% This program can redistributed and/or modified under %%
%% the terms of the LaTeX Project Public License        %%
%% Distributed from CTAN archives in directory          %%
%% macros/latex/base/lppl.txt; either version 1.3c of   %%
%% the License, or (at your option) any later version.  %%
%%
%% DESCRIPTION:
%%   `pst-vehicle' is a PSTricks package
%%
%%
\csname PSTvehicleLoaded\endcsname
\let\PSTvehicleLoaded\endinput
% Requires PSTricks, pst-xkey, pst-node packages
\ifx\PSTricksLoaded\endinput\else\input pstricks.tex\fi
\ifx\PSTXKeyLoaded\endinput\else\input pst-xkey.tex\fi
\ifx\PSTplotLoaded\endinput\else\input pst-plot.tex\fi
\ifx\PSTnodeLoaded\endinput\else\input pst-node.tex\fi
\def\fileversion{1.3}
\def\filedate{2021/12/29}
\message{`PST' v\fileversion, \filedate}

\edef\PstAtCode{\the\catcode`\@} \catcode`\@=11\relax

\definecolor{greenSlp}{rgb}{0,0.4,0.0}
\definecolor{redSlp}{rgb}{0.7,0.134,0.134}

\pst@addfams{pst-vehicle}
\define@key[psset]{pst-vehicle}{backwheel}[\wheelA]{\def\pst@backwheel{#1 }} %
\define@key[psset]{pst-vehicle}{frontwheel}[\wheelA]{\def\pst@frontwheel{#1 }} %
\define@key[psset]{pst-vehicle}{ownvehicle}{\def\pst@ownvehicle{#1 }} %
\define@key[psset]{pst-vehicle}{vehicle}[\Bike]{\def\pst@vehicle{#1}} %
\define@key[psset]{pst-vehicle}{GravNode}[]{\def\pst@GravNode{#1 }} %
\define@key[psset]{pst-vehicle}{d}[5.8]{\def\pst@d{#1 }} %
\define@key[psset]{pst-vehicle}{rF}[1.6]{\def\pst@rF{#1 }} %
\define@key[psset]{pst-vehicle}{rB}[1.6]{\def\pst@rB{#1 }} %
\define@key[psset]{pst-vehicle}{gang}[1]{\def\pst@gang{#1 }} %
\define@key[psset]{pst-vehicle}{epsilon}[1e-6]{\def\pst@epsilon{#1 }} %
\define@key[psset]{pst-vehicle}{startPos}[0]{\def\pst@startPos{#1 }} %
\define@boolkey[psset]{pst-vehicle}[Pst@]{showSlope}[true]{} %
\define@boolkey[psset]{pst-vehicle}[Pst@]{MonoAxis}[false]{} %
%\define@boolkey[psset]{pst-vehicle}[Pst@]{EqValveStartPos}[true]{} %

%% $Id: pst-vehicle.tex 376 2021-12-29 12:28:42Z herbert $
%%
%% This is file `pst-vehicle.tex',
%%
%% IMPORTANT NOTICE:
%%
%% Package `pst-vehicle.tex'
%%
%% Thomas Söll
%% with the collaboration of
%% Juergen Gilg
%% Manuel Luque
%% Herbert Voß (bugfixes)
%%
%% This program can redistributed and/or modified under %%
%% the terms of the LaTeX Project Public License        %%
%% Distributed from CTAN archives in directory          %%
%% macros/latex/base/lppl.txt; either version 1.3c of   %%
%% the License, or (at your option) any later version.  %%
%%
%% DESCRIPTION:
%%   `pst-vehicle' is a PSTricks package
%%
%%
\csname PSTvehicleLoaded\endcsname
\let\PSTvehicleLoaded\endinput
% Requires PSTricks, pst-xkey, pst-node packages
\ifx\PSTricksLoaded\endinput\else\input pstricks.tex\fi
\ifx\PSTXKeyLoaded\endinput\else\input pst-xkey.tex\fi
\ifx\PSTplotLoaded\endinput\else\input pst-plot.tex\fi
\ifx\PSTnodeLoaded\endinput\else\input pst-node.tex\fi
\def\fileversion{1.3}
\def\filedate{2021/12/29}
\message{`PST' v\fileversion, \filedate}

\edef\PstAtCode{\the\catcode`\@} \catcode`\@=11\relax

\definecolor{greenSlp}{rgb}{0,0.4,0.0}
\definecolor{redSlp}{rgb}{0.7,0.134,0.134}

\pst@addfams{pst-vehicle}
\define@key[psset]{pst-vehicle}{backwheel}[\wheelA]{\def\pst@backwheel{#1 }} %
\define@key[psset]{pst-vehicle}{frontwheel}[\wheelA]{\def\pst@frontwheel{#1 }} %
\define@key[psset]{pst-vehicle}{ownvehicle}{\def\pst@ownvehicle{#1 }} %
\define@key[psset]{pst-vehicle}{vehicle}[\Bike]{\def\pst@vehicle{#1}} %
\define@key[psset]{pst-vehicle}{GravNode}[]{\def\pst@GravNode{#1 }} %
\define@key[psset]{pst-vehicle}{d}[5.8]{\def\pst@d{#1 }} %
\define@key[psset]{pst-vehicle}{rF}[1.6]{\def\pst@rF{#1 }} %
\define@key[psset]{pst-vehicle}{rB}[1.6]{\def\pst@rB{#1 }} %
\define@key[psset]{pst-vehicle}{gang}[1]{\def\pst@gang{#1 }} %
\define@key[psset]{pst-vehicle}{epsilon}[1e-6]{\def\pst@epsilon{#1 }} %
\define@key[psset]{pst-vehicle}{startPos}[0]{\def\pst@startPos{#1 }} %
\define@boolkey[psset]{pst-vehicle}[Pst@]{showSlope}[true]{} %
\define@boolkey[psset]{pst-vehicle}[Pst@]{MonoAxis}[false]{} %
%\define@boolkey[psset]{pst-vehicle}[Pst@]{EqValveStartPos}[true]{} %

%% $Id: pst-vehicle.tex 376 2021-12-29 12:28:42Z herbert $
%%
%% This is file `pst-vehicle.tex',
%%
%% IMPORTANT NOTICE:
%%
%% Package `pst-vehicle.tex'
%%
%% Thomas Söll
%% with the collaboration of
%% Juergen Gilg
%% Manuel Luque
%% Herbert Voß (bugfixes)
%%
%% This program can redistributed and/or modified under %%
%% the terms of the LaTeX Project Public License        %%
%% Distributed from CTAN archives in directory          %%
%% macros/latex/base/lppl.txt; either version 1.3c of   %%
%% the License, or (at your option) any later version.  %%
%%
%% DESCRIPTION:
%%   `pst-vehicle' is a PSTricks package
%%
%%
\csname PSTvehicleLoaded\endcsname
\let\PSTvehicleLoaded\endinput
% Requires PSTricks, pst-xkey, pst-node packages
\ifx\PSTricksLoaded\endinput\else\input pstricks.tex\fi
\ifx\PSTXKeyLoaded\endinput\else\input pst-xkey.tex\fi
\ifx\PSTplotLoaded\endinput\else\input pst-plot.tex\fi
\ifx\PSTnodeLoaded\endinput\else\input pst-node.tex\fi
\def\fileversion{1.3}
\def\filedate{2021/12/29}
\message{`PST' v\fileversion, \filedate}

\edef\PstAtCode{\the\catcode`\@} \catcode`\@=11\relax

\definecolor{greenSlp}{rgb}{0,0.4,0.0}
\definecolor{redSlp}{rgb}{0.7,0.134,0.134}

\pst@addfams{pst-vehicle}
\define@key[psset]{pst-vehicle}{backwheel}[\wheelA]{\def\pst@backwheel{#1 }} %
\define@key[psset]{pst-vehicle}{frontwheel}[\wheelA]{\def\pst@frontwheel{#1 }} %
\define@key[psset]{pst-vehicle}{ownvehicle}{\def\pst@ownvehicle{#1 }} %
\define@key[psset]{pst-vehicle}{vehicle}[\Bike]{\def\pst@vehicle{#1}} %
\define@key[psset]{pst-vehicle}{GravNode}[]{\def\pst@GravNode{#1 }} %
\define@key[psset]{pst-vehicle}{d}[5.8]{\def\pst@d{#1 }} %
\define@key[psset]{pst-vehicle}{rF}[1.6]{\def\pst@rF{#1 }} %
\define@key[psset]{pst-vehicle}{rB}[1.6]{\def\pst@rB{#1 }} %
\define@key[psset]{pst-vehicle}{gang}[1]{\def\pst@gang{#1 }} %
\define@key[psset]{pst-vehicle}{epsilon}[1e-6]{\def\pst@epsilon{#1 }} %
\define@key[psset]{pst-vehicle}{startPos}[0]{\def\pst@startPos{#1 }} %
\define@boolkey[psset]{pst-vehicle}[Pst@]{showSlope}[true]{} %
\define@boolkey[psset]{pst-vehicle}[Pst@]{MonoAxis}[false]{} %
%\define@boolkey[psset]{pst-vehicle}[Pst@]{EqValveStartPos}[true]{} %

\input{pst-vehicle.data}

\psset[pst-vehicle]{gang=1,epsilon=1e-6,rB=1.6,rF=1.6,d=5.8,vehicle=\Bike,ownvehicle=,backwheel=\wheelA,frontwheel=\wheelA,showSlope=true,%
startPos=0,MonoAxis=false,GravNode=dA12 2 div 1}
\psset{algebraic}
\def\psVehicle{\def\pst@par{}\pst@object{psVehicle}}%
\def\psVehicle@i#1#2#3{%
\pst@killglue%
\begingroup%
\use@par
%-----------------------------------------------------------------------------------------------------
\expandafter\ifx\pst@vehicle\HighWheeler\psset{rF=4,rB=1.13,d=5.8}\fi
\expandafter\ifx\pst@vehicle\Bike\psset{rF=1.6,rB=1.6,d=5.8}\fi
\expandafter\ifx\pst@vehicle\Truck\psset{rF=1.9,rB=1.9,d=6.28}\fi
\expandafter\ifx\pst@vehicle\Tractor\psset{rF=1,rB=1.4,d=4}\fi
\expandafter\ifx\pst@vehicle\UniCycle\psset{rB=1.6,MonoAxis=true}\fi
\expandafter\ifx\pst@vehicle\Segway\psset{MonoAxis=true}\fi
\begin@SpecialObj
\pst@Verb{%
 /rpn { tx@AlgToPs begin AlgToPs end cvx } def%
 /X0 #2 def % -------------- x-Wert des Punktes auf der Kurve, wo das Hinterrad die Kurve berührt; von diesem Wert startet die Rechnung
 /XST \pst@startPos\space def %   Untergrenze für die Integration bei der Rotationswinkelbestimmung
%---- % Definition of the function f(x), its first derivative f'(x) and  \sqrt{1+f'(x)^2}   Definition de la fonction et la premiere derivee
  /func (#3) rpn def
  /Diff (Derive(1,#3)) rpn def
  /DiffI (Derive(2,#3)) rpn def
  /dAB (sqrt(1+Diff^2)) rpn def
  /dABdiff (Derive(1,sqrt(1+(Derive(1,#3))^2))) rpn def
  /x XST def func /funcxST exch def %---- f(XST)
  /x XST def Diff /DiffxST exch def %---- f'(XST)
  /x X0 def func /funcX0 exch def % ----- f(X0)
  /x X0 def Diff /DiffX0 exch def % ----- f'(X0)
  /x X0 def DiffI /DiffIX0 exch def % --- f''(X0)
%-----------------------------------------------------------------------------------------------------------------------
/eps \pst@epsilon def
% Definition of a transmission between frontwheel and backwheel (interesting for vehicles with pedals) --- Gangschaltung
/Gang \pst@gang def
%% Definition of a scaling factor for the vehicles
#1 /skal exch def
%------------------ Radius frontwheel -----------------------------------------
/rF \pst@rF def
/rB \pst@rB def
/dA12 \pst@d def
/rFs rF skal mul def
/rBs rB skal mul def
%--------% dA12 = Distance between the axes of the wheels ----- Achsabstand --- distance entre les axes
/dA12s dA12 skal mul def
/tA 1 1 DiffX0 dup mul add sqrt div def%
/deltaX0 tA DiffX0 mul rBs mul def
/deltaY0 tA rBs mul def
%-----------------------------------------------------------------------------------------------------------------------------------------
/Function ((x-X0+rBs*DiffX0/(sqrt(1+(DiffX0)^2))-rFs*(Diff)/(sqrt(1+(Diff)^2)))^2+%
            (func-funcX0+rFs/(sqrt(1+(Diff)^2))-rBs/(sqrt(1+(DiffX0)^2)))^2-dA12s^2) rpn def
%-----------------------------------------------------------------------------------------------------------------------------------------
/FunctionST ((x-XST+rBs*DiffxST/(sqrt(1+(DiffxST)^2))-rFs*(Diff)/(sqrt(1+(Diff)^2)))^2+%
            (func-funcxST+rFs/(sqrt(1+(Diff)^2))-rBs/(sqrt(1+(DiffxST)^2)))^2-dA12s^2) rpn def
%--------- inferior value to search for the intersection point of the frontwheel with the curve
/Zeros { %%  Funktion xinf
1 dict begin
/Xinf exch def % ---------------------------------- Untergrenze für die Schnittpunktsuche des Vorderrades mit der Kurve
% superior value to search for the intersection point of the frontwheel with the curve = X0 + distance axes + radius frontwheel + radius backwheel
/Xsup Xinf dA12s add rBs rFs add add def % ---------- Obergrenze für die Schnittpunktsuche des Vorderrades mit der Kurve
%---% Calculating the intersection point of the frontwheel with the function------ Schnittpunktberechnung ---------------------------------
/NB 0 def %-----------------loop-variable  -----  Laufvariable für loop
/NbreIterations 200 def % ---------- Maximum of iterations for the loop
{ %------------------------------------ loop begin ---------------------
 /xM Xinf Xsup add 2 div def %--------- Mittelwert von xM = (Xinf + Xsup):2
   /x Xinf def
   /F_1 FUNK def %----------------- F(Xinf)
   /x xM def %-------------------------
   /F_M FUNK def %----------------- F(xM)
   F_M 0 eq {exit} if %---------------- if F(xM) = 0 --> exit
    F_1 F_M mul 0 ge {/Xinf xM def} %-- F(Xinf) * F(xM) >= 0 Xinf = xM, else Xsup = xM
                     {/Xsup xM def}
    ifelse
 Xinf Xsup sub abs 1e-8 le {exit} if %- if Xinf - Xsup <= 10^-8 --> exit
  /NB NB 1 add def %------------------- else loopvariable NB = NB + 1  Loopvariable um eins erhöhen
 NB NbreIterations ge {exit} if %------ if number of iterations >= 200 --> exit
 } loop
xM
 end
} def
% \ifPst@EqValveStartPos
/FUNK {FunctionST} def XST Zeros /FWxST exch def
% \else /FWxST XST def \fi
/FUNK {Function} def X0 Zeros /FWx exch def
/FWy /x FWx def func def  %----------------------- Berührpunkt des Vorderrades (FWx,FWy)
/mFWy /x FWx def Diff def %----------------------- mFWy = Tangentensteigung in (FWx,FWy)
/TermFW 1 1 mFWy dup mul add sqrt div def %             1/sqrt(1+f'(x_Q)^2)
/deltaxFW TermFW mFWy mul rFs mul def %                  skal*rF*f'(x_Q)*1/sqrt(1+f'(x_Q)^2)
/deltayFW TermFW rFs mul def %                           skal*rF*1/sqrt(1+f'(x_Q)^2)
%------------------------------------------------------------------------------------
/KWRho {DiffI 1 Diff dup mul add 3 exp sqrt div} def
/dPhiB {1 rBs div KWRho sub dAB mul abs} def
/AngleCumB { %
 X1 X2 /x {dPhiB} eps SIMPSON
 } def
 /dPhiF {1 rFs div KWRho sub dAB mul abs} def
/AngleCumF { %
 X1 X2 /x {dPhiF} eps SIMPSON
 } def
% % length of the curve for the wheels  -- La longueur de la courbe pour la roue avant  --- Kurvenlänge
/X1 XST def /X2 X0 def %  Integral_{0}^{X0}
/sB AngleCumB def % ---backwheel - length from 0 to X0 ---- roue arrière ---- Kurvenlänge von 0 bis X0 (Hinterrad)
/X1 FWxST def /X2 FWx def
/sF AngleCumF def % --frontwheel - length from 0 to the abscissa of the intersection frontwheel with curve - roue avant -- Kurvenlänge von 0 bis zum SP Vorderrad - Kurve
%--------------------------------------------------------------------------------------------------------------------
%---% Definition angle of rotation for the backwheel --- Definition de l'angle de roue arriere ---- Rotationswinkel des Hinterrades
/phiB sB RadToDeg neg def
%---% Definition angle of rotation for the frontwheel --- Definition de l'angle de roue avant  ---- Rotationswinkel des Vorderrades
/phiF sF RadToDeg neg def
%---------------------------------------------------------------------------------------------------------------------
%--% coordinates for the axes of the wheels within the non-scaled system -- Koordinaten der Vorderradachse im nicht skalierten System ( außerhalb der \psscalebox )
%----------------------------------------------------------------------------------------------------------------------
/AFx FWx deltaxFW sub def %          % x-coordinate front axis   x-Koordinate der Vorderradachse
/AFy FWy deltayFW add def %          % y-coordinate front axis   y-Koordinate der Vorderradachse
/ABx X0 deltaX0 sub def %            % x-coordinate back axis    x-Koordinate der Hinterradachse
/ABy funcX0 deltaY0 add def %        % y-coordinate back axis    y-Koordinate der Hinterradachse
%--------------- Koordinaten der Vorderradachse im System des Fahrzeugs, also unskalierte Größen verwenden --( innerhalb der \psscalebox )
%--coordinates for the front axis of the wheels within \psscalebox  -- Die Koordinaten der Hinterradachse sind im System des Fahrzeugs  (0,0)
/AF1x dA12 dup mul rF rB sub dup mul sub sqrt def
/AF1y rF rB sub def
%------% slope and angle of the vehicle-----------------------------------------------------------------------------
/m-vehicle AFy ABy sub AFx ABx sub div def
/beta m-vehicle 1 atan def %          % angle for both radii of front- and backwheel are equal -- Neigungswinkel des Fahrzeugs bei gleich großen Rädern
/alpha AF1y neg AF1x atan def %       % additional angle if radii of front- and backwheel are not equal --  zusätzlicher Winkel bei unterschiedlich großen Rädern
% whole angle (correction with +180 in case the whole angle gets 90 degrees which can be possible with not equal radii)
/gamma beta alpha add AFx ABx lt { 180 add } if def %   gesamter Neigungswinkel des Fahrzeugs
%--------% Special case (mono-axis vehicle) Nr 4 --> segway ---------------------------------------------------------------------------
\ifPst@MonoAxis DiffX0 1 atan /omega exch def DiffX0 /mVehicle exch def \else FWy funcX0 sub FWx X0 sub div /mVehicle exch def mVehicle 1 atan /omega exch def \fi
/normNorm FWy funcX0 eq { 1 } { X0 FWx sub FWy funcX0 sub div dup mul 1 add sqrt } ifelse def
/mTgy FWy funcX0 eq { 1 } { X0 FWx sub FWy funcX0 sub div normNorm div } ifelse def
/mTgx FWy funcX0 eq { 0 } { 1 normNorm div } ifelse def
/xMTg X0 FWx add 2 div def
/yMTg funcX0 FWy add 2 div def
}%
%------------------------%%%%  END OF PS-CODE %%%%%--------------------------------------------------------
\pnode(!FWx FWy){FW}%
% % angle between the axes   angle de droite entre les axes
\rput{!gamma}(!ABx ABy){%  Das Fahrzeug wird mit Hilfe der Hinterradachse (ABx,ABy) und Gesamtdrehwinkel gamma gesetzt
%------------------------- % SETTING SOME VEHICLES  -------------------------------------------------------
\psscalebox{#1}{%---------------Das Fahrzeug kann skaliert werden -----------------------------------------
%\psset{linecolor=#2}%
\pnode(!\pst@GravNode){GravC}%
\pst@vehicle}}%
\ifPst@showSlope
\rput{!omega}(!X0 mVehicle 0 ge { 0 add } { 100 add } ifelse funcX0 0.75 sub){%
\psframebox[fillstyle=solid,fillcolor=greenSlp,linestyle=none]{\footnotesize\color{white}$m\geq 0$}}
\rput{!omega}(!X0 mVehicle 0 lt { 0 add } { 100 add } ifelse funcX0 0.75 sub){%
\psframebox[fillstyle=solid,fillcolor=redSlp,linestyle=none]{\footnotesize\color{white}$m < 0$}}
\ifPst@MonoAxis
\psplotTangent[linewidth=0.75\pslinewidth,linecolor=orange,nodesep=-2]{X0}{1}{#3}
\psplotTangent[linewidth=0.75\pslinewidth,linecolor=greenSlp,Tnormal,nodesep=-2]{X0}{1}{#3}
\psdot[linecolor=red](!X0 funcX0)
\else
\pcline[linecolor=greenSlp,nodesepA=-1,nodesepB=-3](!xMTg yMTg)(!xMTg mTgy 0 ge { mTgx } { mTgx neg } ifelse add yMTg mTgy abs add)
\pcline[linecolor=magenta,nodesep=-2](FW)(!X0 funcX0)
\psdot[linecolor=red](FW)
\psdot[linecolor=red](!X0 funcX0)
\fi%
\fi%
\showpointsfalse
\end@SpecialObj
\endgroup\ignorespaces}%
%
\def\SlopeoMeter#1#2{%
\colorlet{slpmColor}{#1}
\rput{0}(0,0){%
\pscircle[fillstyle=solid,fillcolor=black!90,linewidth=0.5pt,linestyle=solid](0,0){2.2}
\rput{!#2}(0,0){%
\multido{\r=.500+-.008,\rC=0+0.02}{25}{%
\pscircle[linewidth=0.008,linecolor=slpmColor,strokeopacity=\rC](0,0){!\r}
}
\multido{\i=0+2,\r=.400+-.008,\rC=0+0.02}{50}{%
\pswedge[linestyle=none,linewidth=0,fillstyle=solid,fillcolor=slpmColor,opacity=\r](0,0){2}{\i}{!\i\space 2 add}
}}
\multido{\nA=-90+10}{19}{%
\psline[linecolor=slpmColor](1.8;\nA)(1.95;\nA)
\rput(1.65;\nA){\psscalebox{0.4}{\color{slpmColor!10}\nA}}
}
\rput(-0.95,0.85){\psscalebox{0.6}{\color{slpmColor!20}\texttt{Steigungswinkel}}}
\rput(-0.95,0.5){\psscalebox{0.75}{\color{slpmColor!40}\texttt{Slope-o-Meter}}}
\rput{!#2}(0,0){%
\pspolygon[fillstyle=solid,fillcolor=slpmColor!50,linecolor=slpmColor,linejoin=2,linewidth=0.1pt](0.1;15)(0.1;-15)(1.75;0)
\pscircle[fillstyle=solid,fillcolor=black,linestyle=none,linewidth=0pt](0,0){0.3}
}}}
\catcode`\@=\PstAtCode\relax
\endinput


\psset[pst-vehicle]{gang=1,epsilon=1e-6,rB=1.6,rF=1.6,d=5.8,vehicle=\Bike,ownvehicle=,backwheel=\wheelA,frontwheel=\wheelA,showSlope=true,%
startPos=0,MonoAxis=false,GravNode=dA12 2 div 1}
\psset{algebraic}
\def\psVehicle{\def\pst@par{}\pst@object{psVehicle}}%
\def\psVehicle@i#1#2#3{%
\pst@killglue%
\begingroup%
\use@par
%-----------------------------------------------------------------------------------------------------
\expandafter\ifx\pst@vehicle\HighWheeler\psset{rF=4,rB=1.13,d=5.8}\fi
\expandafter\ifx\pst@vehicle\Bike\psset{rF=1.6,rB=1.6,d=5.8}\fi
\expandafter\ifx\pst@vehicle\Truck\psset{rF=1.9,rB=1.9,d=6.28}\fi
\expandafter\ifx\pst@vehicle\Tractor\psset{rF=1,rB=1.4,d=4}\fi
\expandafter\ifx\pst@vehicle\UniCycle\psset{rB=1.6,MonoAxis=true}\fi
\expandafter\ifx\pst@vehicle\Segway\psset{MonoAxis=true}\fi
\begin@SpecialObj
\pst@Verb{%
 /rpn { tx@AlgToPs begin AlgToPs end cvx } def%
 /X0 #2 def % -------------- x-Wert des Punktes auf der Kurve, wo das Hinterrad die Kurve berührt; von diesem Wert startet die Rechnung
 /XST \pst@startPos\space def %   Untergrenze für die Integration bei der Rotationswinkelbestimmung
%---- % Definition of the function f(x), its first derivative f'(x) and  \sqrt{1+f'(x)^2}   Definition de la fonction et la premiere derivee
  /func (#3) rpn def
  /Diff (Derive(1,#3)) rpn def
  /DiffI (Derive(2,#3)) rpn def
  /dAB (sqrt(1+Diff^2)) rpn def
  /dABdiff (Derive(1,sqrt(1+(Derive(1,#3))^2))) rpn def
  /x XST def func /funcxST exch def %---- f(XST)
  /x XST def Diff /DiffxST exch def %---- f'(XST)
  /x X0 def func /funcX0 exch def % ----- f(X0)
  /x X0 def Diff /DiffX0 exch def % ----- f'(X0)
  /x X0 def DiffI /DiffIX0 exch def % --- f''(X0)
%-----------------------------------------------------------------------------------------------------------------------
/eps \pst@epsilon def
% Definition of a transmission between frontwheel and backwheel (interesting for vehicles with pedals) --- Gangschaltung
/Gang \pst@gang def
%% Definition of a scaling factor for the vehicles
#1 /skal exch def
%------------------ Radius frontwheel -----------------------------------------
/rF \pst@rF def
/rB \pst@rB def
/dA12 \pst@d def
/rFs rF skal mul def
/rBs rB skal mul def
%--------% dA12 = Distance between the axes of the wheels ----- Achsabstand --- distance entre les axes
/dA12s dA12 skal mul def
/tA 1 1 DiffX0 dup mul add sqrt div def%
/deltaX0 tA DiffX0 mul rBs mul def
/deltaY0 tA rBs mul def
%-----------------------------------------------------------------------------------------------------------------------------------------
/Function ((x-X0+rBs*DiffX0/(sqrt(1+(DiffX0)^2))-rFs*(Diff)/(sqrt(1+(Diff)^2)))^2+%
            (func-funcX0+rFs/(sqrt(1+(Diff)^2))-rBs/(sqrt(1+(DiffX0)^2)))^2-dA12s^2) rpn def
%-----------------------------------------------------------------------------------------------------------------------------------------
/FunctionST ((x-XST+rBs*DiffxST/(sqrt(1+(DiffxST)^2))-rFs*(Diff)/(sqrt(1+(Diff)^2)))^2+%
            (func-funcxST+rFs/(sqrt(1+(Diff)^2))-rBs/(sqrt(1+(DiffxST)^2)))^2-dA12s^2) rpn def
%--------- inferior value to search for the intersection point of the frontwheel with the curve
/Zeros { %%  Funktion xinf
1 dict begin
/Xinf exch def % ---------------------------------- Untergrenze für die Schnittpunktsuche des Vorderrades mit der Kurve
% superior value to search for the intersection point of the frontwheel with the curve = X0 + distance axes + radius frontwheel + radius backwheel
/Xsup Xinf dA12s add rBs rFs add add def % ---------- Obergrenze für die Schnittpunktsuche des Vorderrades mit der Kurve
%---% Calculating the intersection point of the frontwheel with the function------ Schnittpunktberechnung ---------------------------------
/NB 0 def %-----------------loop-variable  -----  Laufvariable für loop
/NbreIterations 200 def % ---------- Maximum of iterations for the loop
{ %------------------------------------ loop begin ---------------------
 /xM Xinf Xsup add 2 div def %--------- Mittelwert von xM = (Xinf + Xsup):2
   /x Xinf def
   /F_1 FUNK def %----------------- F(Xinf)
   /x xM def %-------------------------
   /F_M FUNK def %----------------- F(xM)
   F_M 0 eq {exit} if %---------------- if F(xM) = 0 --> exit
    F_1 F_M mul 0 ge {/Xinf xM def} %-- F(Xinf) * F(xM) >= 0 Xinf = xM, else Xsup = xM
                     {/Xsup xM def}
    ifelse
 Xinf Xsup sub abs 1e-8 le {exit} if %- if Xinf - Xsup <= 10^-8 --> exit
  /NB NB 1 add def %------------------- else loopvariable NB = NB + 1  Loopvariable um eins erhöhen
 NB NbreIterations ge {exit} if %------ if number of iterations >= 200 --> exit
 } loop
xM
 end
} def
% \ifPst@EqValveStartPos
/FUNK {FunctionST} def XST Zeros /FWxST exch def
% \else /FWxST XST def \fi
/FUNK {Function} def X0 Zeros /FWx exch def
/FWy /x FWx def func def  %----------------------- Berührpunkt des Vorderrades (FWx,FWy)
/mFWy /x FWx def Diff def %----------------------- mFWy = Tangentensteigung in (FWx,FWy)
/TermFW 1 1 mFWy dup mul add sqrt div def %             1/sqrt(1+f'(x_Q)^2)
/deltaxFW TermFW mFWy mul rFs mul def %                  skal*rF*f'(x_Q)*1/sqrt(1+f'(x_Q)^2)
/deltayFW TermFW rFs mul def %                           skal*rF*1/sqrt(1+f'(x_Q)^2)
%------------------------------------------------------------------------------------
/KWRho {DiffI 1 Diff dup mul add 3 exp sqrt div} def
/dPhiB {1 rBs div KWRho sub dAB mul abs} def
/AngleCumB { %
 X1 X2 /x {dPhiB} eps SIMPSON
 } def
 /dPhiF {1 rFs div KWRho sub dAB mul abs} def
/AngleCumF { %
 X1 X2 /x {dPhiF} eps SIMPSON
 } def
% % length of the curve for the wheels  -- La longueur de la courbe pour la roue avant  --- Kurvenlänge
/X1 XST def /X2 X0 def %  Integral_{0}^{X0}
/sB AngleCumB def % ---backwheel - length from 0 to X0 ---- roue arrière ---- Kurvenlänge von 0 bis X0 (Hinterrad)
/X1 FWxST def /X2 FWx def
/sF AngleCumF def % --frontwheel - length from 0 to the abscissa of the intersection frontwheel with curve - roue avant -- Kurvenlänge von 0 bis zum SP Vorderrad - Kurve
%--------------------------------------------------------------------------------------------------------------------
%---% Definition angle of rotation for the backwheel --- Definition de l'angle de roue arriere ---- Rotationswinkel des Hinterrades
/phiB sB RadToDeg neg def
%---% Definition angle of rotation for the frontwheel --- Definition de l'angle de roue avant  ---- Rotationswinkel des Vorderrades
/phiF sF RadToDeg neg def
%---------------------------------------------------------------------------------------------------------------------
%--% coordinates for the axes of the wheels within the non-scaled system -- Koordinaten der Vorderradachse im nicht skalierten System ( außerhalb der \psscalebox )
%----------------------------------------------------------------------------------------------------------------------
/AFx FWx deltaxFW sub def %          % x-coordinate front axis   x-Koordinate der Vorderradachse
/AFy FWy deltayFW add def %          % y-coordinate front axis   y-Koordinate der Vorderradachse
/ABx X0 deltaX0 sub def %            % x-coordinate back axis    x-Koordinate der Hinterradachse
/ABy funcX0 deltaY0 add def %        % y-coordinate back axis    y-Koordinate der Hinterradachse
%--------------- Koordinaten der Vorderradachse im System des Fahrzeugs, also unskalierte Größen verwenden --( innerhalb der \psscalebox )
%--coordinates for the front axis of the wheels within \psscalebox  -- Die Koordinaten der Hinterradachse sind im System des Fahrzeugs  (0,0)
/AF1x dA12 dup mul rF rB sub dup mul sub sqrt def
/AF1y rF rB sub def
%------% slope and angle of the vehicle-----------------------------------------------------------------------------
/m-vehicle AFy ABy sub AFx ABx sub div def
/beta m-vehicle 1 atan def %          % angle for both radii of front- and backwheel are equal -- Neigungswinkel des Fahrzeugs bei gleich großen Rädern
/alpha AF1y neg AF1x atan def %       % additional angle if radii of front- and backwheel are not equal --  zusätzlicher Winkel bei unterschiedlich großen Rädern
% whole angle (correction with +180 in case the whole angle gets 90 degrees which can be possible with not equal radii)
/gamma beta alpha add AFx ABx lt { 180 add } if def %   gesamter Neigungswinkel des Fahrzeugs
%--------% Special case (mono-axis vehicle) Nr 4 --> segway ---------------------------------------------------------------------------
\ifPst@MonoAxis DiffX0 1 atan /omega exch def DiffX0 /mVehicle exch def \else FWy funcX0 sub FWx X0 sub div /mVehicle exch def mVehicle 1 atan /omega exch def \fi
/normNorm FWy funcX0 eq { 1 } { X0 FWx sub FWy funcX0 sub div dup mul 1 add sqrt } ifelse def
/mTgy FWy funcX0 eq { 1 } { X0 FWx sub FWy funcX0 sub div normNorm div } ifelse def
/mTgx FWy funcX0 eq { 0 } { 1 normNorm div } ifelse def
/xMTg X0 FWx add 2 div def
/yMTg funcX0 FWy add 2 div def
}%
%------------------------%%%%  END OF PS-CODE %%%%%--------------------------------------------------------
\pnode(!FWx FWy){FW}%
% % angle between the axes   angle de droite entre les axes
\rput{!gamma}(!ABx ABy){%  Das Fahrzeug wird mit Hilfe der Hinterradachse (ABx,ABy) und Gesamtdrehwinkel gamma gesetzt
%------------------------- % SETTING SOME VEHICLES  -------------------------------------------------------
\psscalebox{#1}{%---------------Das Fahrzeug kann skaliert werden -----------------------------------------
%\psset{linecolor=#2}%
\pnode(!\pst@GravNode){GravC}%
\pst@vehicle}}%
\ifPst@showSlope
\rput{!omega}(!X0 mVehicle 0 ge { 0 add } { 100 add } ifelse funcX0 0.75 sub){%
\psframebox[fillstyle=solid,fillcolor=greenSlp,linestyle=none]{\footnotesize\color{white}$m\geq 0$}}
\rput{!omega}(!X0 mVehicle 0 lt { 0 add } { 100 add } ifelse funcX0 0.75 sub){%
\psframebox[fillstyle=solid,fillcolor=redSlp,linestyle=none]{\footnotesize\color{white}$m < 0$}}
\ifPst@MonoAxis
\psplotTangent[linewidth=0.75\pslinewidth,linecolor=orange,nodesep=-2]{X0}{1}{#3}
\psplotTangent[linewidth=0.75\pslinewidth,linecolor=greenSlp,Tnormal,nodesep=-2]{X0}{1}{#3}
\psdot[linecolor=red](!X0 funcX0)
\else
\pcline[linecolor=greenSlp,nodesepA=-1,nodesepB=-3](!xMTg yMTg)(!xMTg mTgy 0 ge { mTgx } { mTgx neg } ifelse add yMTg mTgy abs add)
\pcline[linecolor=magenta,nodesep=-2](FW)(!X0 funcX0)
\psdot[linecolor=red](FW)
\psdot[linecolor=red](!X0 funcX0)
\fi%
\fi%
\showpointsfalse
\end@SpecialObj
\endgroup\ignorespaces}%
%
\def\SlopeoMeter#1#2{%
\colorlet{slpmColor}{#1}
\rput{0}(0,0){%
\pscircle[fillstyle=solid,fillcolor=black!90,linewidth=0.5pt,linestyle=solid](0,0){2.2}
\rput{!#2}(0,0){%
\multido{\r=.500+-.008,\rC=0+0.02}{25}{%
\pscircle[linewidth=0.008,linecolor=slpmColor,strokeopacity=\rC](0,0){!\r}
}
\multido{\i=0+2,\r=.400+-.008,\rC=0+0.02}{50}{%
\pswedge[linestyle=none,linewidth=0,fillstyle=solid,fillcolor=slpmColor,opacity=\r](0,0){2}{\i}{!\i\space 2 add}
}}
\multido{\nA=-90+10}{19}{%
\psline[linecolor=slpmColor](1.8;\nA)(1.95;\nA)
\rput(1.65;\nA){\psscalebox{0.4}{\color{slpmColor!10}\nA}}
}
\rput(-0.95,0.85){\psscalebox{0.6}{\color{slpmColor!20}\texttt{Steigungswinkel}}}
\rput(-0.95,0.5){\psscalebox{0.75}{\color{slpmColor!40}\texttt{Slope-o-Meter}}}
\rput{!#2}(0,0){%
\pspolygon[fillstyle=solid,fillcolor=slpmColor!50,linecolor=slpmColor,linejoin=2,linewidth=0.1pt](0.1;15)(0.1;-15)(1.75;0)
\pscircle[fillstyle=solid,fillcolor=black,linestyle=none,linewidth=0pt](0,0){0.3}
}}}
\catcode`\@=\PstAtCode\relax
\endinput


\psset[pst-vehicle]{gang=1,epsilon=1e-6,rB=1.6,rF=1.6,d=5.8,vehicle=\Bike,ownvehicle=,backwheel=\wheelA,frontwheel=\wheelA,showSlope=true,%
startPos=0,MonoAxis=false,GravNode=dA12 2 div 1}
\psset{algebraic}
\def\psVehicle{\def\pst@par{}\pst@object{psVehicle}}%
\def\psVehicle@i#1#2#3{%
\pst@killglue%
\begingroup%
\use@par
%-----------------------------------------------------------------------------------------------------
\expandafter\ifx\pst@vehicle\HighWheeler\psset{rF=4,rB=1.13,d=5.8}\fi
\expandafter\ifx\pst@vehicle\Bike\psset{rF=1.6,rB=1.6,d=5.8}\fi
\expandafter\ifx\pst@vehicle\Truck\psset{rF=1.9,rB=1.9,d=6.28}\fi
\expandafter\ifx\pst@vehicle\Tractor\psset{rF=1,rB=1.4,d=4}\fi
\expandafter\ifx\pst@vehicle\UniCycle\psset{rB=1.6,MonoAxis=true}\fi
\expandafter\ifx\pst@vehicle\Segway\psset{MonoAxis=true}\fi
\begin@SpecialObj
\pst@Verb{%
 /rpn { tx@AlgToPs begin AlgToPs end cvx } def%
 /X0 #2 def % -------------- x-Wert des Punktes auf der Kurve, wo das Hinterrad die Kurve berührt; von diesem Wert startet die Rechnung
 /XST \pst@startPos\space def %   Untergrenze für die Integration bei der Rotationswinkelbestimmung
%---- % Definition of the function f(x), its first derivative f'(x) and  \sqrt{1+f'(x)^2}   Definition de la fonction et la premiere derivee
  /func (#3) rpn def
  /Diff (Derive(1,#3)) rpn def
  /DiffI (Derive(2,#3)) rpn def
  /dAB (sqrt(1+Diff^2)) rpn def
  /dABdiff (Derive(1,sqrt(1+(Derive(1,#3))^2))) rpn def
  /x XST def func /funcxST exch def %---- f(XST)
  /x XST def Diff /DiffxST exch def %---- f'(XST)
  /x X0 def func /funcX0 exch def % ----- f(X0)
  /x X0 def Diff /DiffX0 exch def % ----- f'(X0)
  /x X0 def DiffI /DiffIX0 exch def % --- f''(X0)
%-----------------------------------------------------------------------------------------------------------------------
/eps \pst@epsilon def
% Definition of a transmission between frontwheel and backwheel (interesting for vehicles with pedals) --- Gangschaltung
/Gang \pst@gang def
%% Definition of a scaling factor for the vehicles
#1 /skal exch def
%------------------ Radius frontwheel -----------------------------------------
/rF \pst@rF def
/rB \pst@rB def
/dA12 \pst@d def
/rFs rF skal mul def
/rBs rB skal mul def
%--------% dA12 = Distance between the axes of the wheels ----- Achsabstand --- distance entre les axes
/dA12s dA12 skal mul def
/tA 1 1 DiffX0 dup mul add sqrt div def%
/deltaX0 tA DiffX0 mul rBs mul def
/deltaY0 tA rBs mul def
%-----------------------------------------------------------------------------------------------------------------------------------------
/Function ((x-X0+rBs*DiffX0/(sqrt(1+(DiffX0)^2))-rFs*(Diff)/(sqrt(1+(Diff)^2)))^2+%
            (func-funcX0+rFs/(sqrt(1+(Diff)^2))-rBs/(sqrt(1+(DiffX0)^2)))^2-dA12s^2) rpn def
%-----------------------------------------------------------------------------------------------------------------------------------------
/FunctionST ((x-XST+rBs*DiffxST/(sqrt(1+(DiffxST)^2))-rFs*(Diff)/(sqrt(1+(Diff)^2)))^2+%
            (func-funcxST+rFs/(sqrt(1+(Diff)^2))-rBs/(sqrt(1+(DiffxST)^2)))^2-dA12s^2) rpn def
%--------- inferior value to search for the intersection point of the frontwheel with the curve
/Zeros { %%  Funktion xinf
1 dict begin
/Xinf exch def % ---------------------------------- Untergrenze für die Schnittpunktsuche des Vorderrades mit der Kurve
% superior value to search for the intersection point of the frontwheel with the curve = X0 + distance axes + radius frontwheel + radius backwheel
/Xsup Xinf dA12s add rBs rFs add add def % ---------- Obergrenze für die Schnittpunktsuche des Vorderrades mit der Kurve
%---% Calculating the intersection point of the frontwheel with the function------ Schnittpunktberechnung ---------------------------------
/NB 0 def %-----------------loop-variable  -----  Laufvariable für loop
/NbreIterations 200 def % ---------- Maximum of iterations for the loop
{ %------------------------------------ loop begin ---------------------
 /xM Xinf Xsup add 2 div def %--------- Mittelwert von xM = (Xinf + Xsup):2
   /x Xinf def
   /F_1 FUNK def %----------------- F(Xinf)
   /x xM def %-------------------------
   /F_M FUNK def %----------------- F(xM)
   F_M 0 eq {exit} if %---------------- if F(xM) = 0 --> exit
    F_1 F_M mul 0 ge {/Xinf xM def} %-- F(Xinf) * F(xM) >= 0 Xinf = xM, else Xsup = xM
                     {/Xsup xM def}
    ifelse
 Xinf Xsup sub abs 1e-8 le {exit} if %- if Xinf - Xsup <= 10^-8 --> exit
  /NB NB 1 add def %------------------- else loopvariable NB = NB + 1  Loopvariable um eins erhöhen
 NB NbreIterations ge {exit} if %------ if number of iterations >= 200 --> exit
 } loop
xM
 end
} def
% \ifPst@EqValveStartPos
/FUNK {FunctionST} def XST Zeros /FWxST exch def
% \else /FWxST XST def \fi
/FUNK {Function} def X0 Zeros /FWx exch def
/FWy /x FWx def func def  %----------------------- Berührpunkt des Vorderrades (FWx,FWy)
/mFWy /x FWx def Diff def %----------------------- mFWy = Tangentensteigung in (FWx,FWy)
/TermFW 1 1 mFWy dup mul add sqrt div def %             1/sqrt(1+f'(x_Q)^2)
/deltaxFW TermFW mFWy mul rFs mul def %                  skal*rF*f'(x_Q)*1/sqrt(1+f'(x_Q)^2)
/deltayFW TermFW rFs mul def %                           skal*rF*1/sqrt(1+f'(x_Q)^2)
%------------------------------------------------------------------------------------
/KWRho {DiffI 1 Diff dup mul add 3 exp sqrt div} def
/dPhiB {1 rBs div KWRho sub dAB mul abs} def
/AngleCumB { %
 X1 X2 /x {dPhiB} eps SIMPSON
 } def
 /dPhiF {1 rFs div KWRho sub dAB mul abs} def
/AngleCumF { %
 X1 X2 /x {dPhiF} eps SIMPSON
 } def
% % length of the curve for the wheels  -- La longueur de la courbe pour la roue avant  --- Kurvenlänge
/X1 XST def /X2 X0 def %  Integral_{0}^{X0}
/sB AngleCumB def % ---backwheel - length from 0 to X0 ---- roue arrière ---- Kurvenlänge von 0 bis X0 (Hinterrad)
/X1 FWxST def /X2 FWx def
/sF AngleCumF def % --frontwheel - length from 0 to the abscissa of the intersection frontwheel with curve - roue avant -- Kurvenlänge von 0 bis zum SP Vorderrad - Kurve
%--------------------------------------------------------------------------------------------------------------------
%---% Definition angle of rotation for the backwheel --- Definition de l'angle de roue arriere ---- Rotationswinkel des Hinterrades
/phiB sB RadToDeg neg def
%---% Definition angle of rotation for the frontwheel --- Definition de l'angle de roue avant  ---- Rotationswinkel des Vorderrades
/phiF sF RadToDeg neg def
%---------------------------------------------------------------------------------------------------------------------
%--% coordinates for the axes of the wheels within the non-scaled system -- Koordinaten der Vorderradachse im nicht skalierten System ( außerhalb der \psscalebox )
%----------------------------------------------------------------------------------------------------------------------
/AFx FWx deltaxFW sub def %          % x-coordinate front axis   x-Koordinate der Vorderradachse
/AFy FWy deltayFW add def %          % y-coordinate front axis   y-Koordinate der Vorderradachse
/ABx X0 deltaX0 sub def %            % x-coordinate back axis    x-Koordinate der Hinterradachse
/ABy funcX0 deltaY0 add def %        % y-coordinate back axis    y-Koordinate der Hinterradachse
%--------------- Koordinaten der Vorderradachse im System des Fahrzeugs, also unskalierte Größen verwenden --( innerhalb der \psscalebox )
%--coordinates for the front axis of the wheels within \psscalebox  -- Die Koordinaten der Hinterradachse sind im System des Fahrzeugs  (0,0)
/AF1x dA12 dup mul rF rB sub dup mul sub sqrt def
/AF1y rF rB sub def
%------% slope and angle of the vehicle-----------------------------------------------------------------------------
/m-vehicle AFy ABy sub AFx ABx sub div def
/beta m-vehicle 1 atan def %          % angle for both radii of front- and backwheel are equal -- Neigungswinkel des Fahrzeugs bei gleich großen Rädern
/alpha AF1y neg AF1x atan def %       % additional angle if radii of front- and backwheel are not equal --  zusätzlicher Winkel bei unterschiedlich großen Rädern
% whole angle (correction with +180 in case the whole angle gets 90 degrees which can be possible with not equal radii)
/gamma beta alpha add AFx ABx lt { 180 add } if def %   gesamter Neigungswinkel des Fahrzeugs
%--------% Special case (mono-axis vehicle) Nr 4 --> segway ---------------------------------------------------------------------------
\ifPst@MonoAxis DiffX0 1 atan /omega exch def DiffX0 /mVehicle exch def \else FWy funcX0 sub FWx X0 sub div /mVehicle exch def mVehicle 1 atan /omega exch def \fi
/normNorm FWy funcX0 eq { 1 } { X0 FWx sub FWy funcX0 sub div dup mul 1 add sqrt } ifelse def
/mTgy FWy funcX0 eq { 1 } { X0 FWx sub FWy funcX0 sub div normNorm div } ifelse def
/mTgx FWy funcX0 eq { 0 } { 1 normNorm div } ifelse def
/xMTg X0 FWx add 2 div def
/yMTg funcX0 FWy add 2 div def
}%
%------------------------%%%%  END OF PS-CODE %%%%%--------------------------------------------------------
\pnode(!FWx FWy){FW}%
% % angle between the axes   angle de droite entre les axes
\rput{!gamma}(!ABx ABy){%  Das Fahrzeug wird mit Hilfe der Hinterradachse (ABx,ABy) und Gesamtdrehwinkel gamma gesetzt
%------------------------- % SETTING SOME VEHICLES  -------------------------------------------------------
\psscalebox{#1}{%---------------Das Fahrzeug kann skaliert werden -----------------------------------------
%\psset{linecolor=#2}%
\pnode(!\pst@GravNode){GravC}%
\pst@vehicle}}%
\ifPst@showSlope
\rput{!omega}(!X0 mVehicle 0 ge { 0 add } { 100 add } ifelse funcX0 0.75 sub){%
\psframebox[fillstyle=solid,fillcolor=greenSlp,linestyle=none]{\footnotesize\color{white}$m\geq 0$}}
\rput{!omega}(!X0 mVehicle 0 lt { 0 add } { 100 add } ifelse funcX0 0.75 sub){%
\psframebox[fillstyle=solid,fillcolor=redSlp,linestyle=none]{\footnotesize\color{white}$m < 0$}}
\ifPst@MonoAxis
\psplotTangent[linewidth=0.75\pslinewidth,linecolor=orange,nodesep=-2]{X0}{1}{#3}
\psplotTangent[linewidth=0.75\pslinewidth,linecolor=greenSlp,Tnormal,nodesep=-2]{X0}{1}{#3}
\psdot[linecolor=red](!X0 funcX0)
\else
\pcline[linecolor=greenSlp,nodesepA=-1,nodesepB=-3](!xMTg yMTg)(!xMTg mTgy 0 ge { mTgx } { mTgx neg } ifelse add yMTg mTgy abs add)
\pcline[linecolor=magenta,nodesep=-2](FW)(!X0 funcX0)
\psdot[linecolor=red](FW)
\psdot[linecolor=red](!X0 funcX0)
\fi%
\fi%
\showpointsfalse
\end@SpecialObj
\endgroup\ignorespaces}%
%
\def\SlopeoMeter#1#2{%
\colorlet{slpmColor}{#1}
\rput{0}(0,0){%
\pscircle[fillstyle=solid,fillcolor=black!90,linewidth=0.5pt,linestyle=solid](0,0){2.2}
\rput{!#2}(0,0){%
\multido{\r=.500+-.008,\rC=0+0.02}{25}{%
\pscircle[linewidth=0.008,linecolor=slpmColor,strokeopacity=\rC](0,0){!\r}
}
\multido{\i=0+2,\r=.400+-.008,\rC=0+0.02}{50}{%
\pswedge[linestyle=none,linewidth=0,fillstyle=solid,fillcolor=slpmColor,opacity=\r](0,0){2}{\i}{!\i\space 2 add}
}}
\multido{\nA=-90+10}{19}{%
\psline[linecolor=slpmColor](1.8;\nA)(1.95;\nA)
\rput(1.65;\nA){\psscalebox{0.4}{\color{slpmColor!10}\nA}}
}
\rput(-0.95,0.85){\psscalebox{0.6}{\color{slpmColor!20}\texttt{Steigungswinkel}}}
\rput(-0.95,0.5){\psscalebox{0.75}{\color{slpmColor!40}\texttt{Slope-o-Meter}}}
\rput{!#2}(0,0){%
\pspolygon[fillstyle=solid,fillcolor=slpmColor!50,linecolor=slpmColor,linejoin=2,linewidth=0.1pt](0.1;15)(0.1;-15)(1.75;0)
\pscircle[fillstyle=solid,fillcolor=black,linestyle=none,linewidth=0pt](0,0){0.3}
}}}
\catcode`\@=\PstAtCode\relax
\endinput


\psset[pst-vehicle]{gang=1,epsilon=1e-6,rB=1.6,rF=1.6,d=5.8,vehicle=\Bike,ownvehicle=,backwheel=\wheelA,frontwheel=\wheelA,showSlope=true,%
startPos=0,MonoAxis=false,GravNode=dA12 2 div 1}
\psset{algebraic}
\def\psVehicle{\def\pst@par{}\pst@object{psVehicle}}%
\def\psVehicle@i#1#2#3{%
\pst@killglue%
\begingroup%
\use@par
%-----------------------------------------------------------------------------------------------------
\expandafter\ifx\pst@vehicle\HighWheeler\psset{rF=4,rB=1.13,d=5.8}\fi
\expandafter\ifx\pst@vehicle\Bike\psset{rF=1.6,rB=1.6,d=5.8}\fi
\expandafter\ifx\pst@vehicle\Truck\psset{rF=1.9,rB=1.9,d=6.28}\fi
\expandafter\ifx\pst@vehicle\Tractor\psset{rF=1,rB=1.4,d=4}\fi
\expandafter\ifx\pst@vehicle\UniCycle\psset{rB=1.6,MonoAxis=true}\fi
\expandafter\ifx\pst@vehicle\Segway\psset{MonoAxis=true}\fi
\begin@SpecialObj
\pst@Verb{%
 /rpn { tx@AlgToPs begin AlgToPs end cvx } def%
 /X0 #2 def % -------------- x-Wert des Punktes auf der Kurve, wo das Hinterrad die Kurve berührt; von diesem Wert startet die Rechnung
 /XST \pst@startPos\space def %   Untergrenze für die Integration bei der Rotationswinkelbestimmung
%---- % Definition of the function f(x), its first derivative f'(x) and  \sqrt{1+f'(x)^2}   Definition de la fonction et la premiere derivee
  /func (#3) rpn def
  /Diff (Derive(1,#3)) rpn def
  /DiffI (Derive(2,#3)) rpn def
  /dAB (sqrt(1+Diff^2)) rpn def
  /dABdiff (Derive(1,sqrt(1+(Derive(1,#3))^2))) rpn def
  /x XST def func /funcxST exch def %---- f(XST)
  /x XST def Diff /DiffxST exch def %---- f'(XST)
  /x X0 def func /funcX0 exch def % ----- f(X0)
  /x X0 def Diff /DiffX0 exch def % ----- f'(X0)
  /x X0 def DiffI /DiffIX0 exch def % --- f''(X0)
%-----------------------------------------------------------------------------------------------------------------------
/eps \pst@epsilon def
% Definition of a transmission between frontwheel and backwheel (interesting for vehicles with pedals) --- Gangschaltung
/Gang \pst@gang def
%% Definition of a scaling factor for the vehicles
#1 /skal exch def
%------------------ Radius frontwheel -----------------------------------------
/rF \pst@rF def
/rB \pst@rB def
/dA12 \pst@d def
/rFs rF skal mul def
/rBs rB skal mul def
%--------% dA12 = Distance between the axes of the wheels ----- Achsabstand --- distance entre les axes
/dA12s dA12 skal mul def
/tA 1 1 DiffX0 dup mul add sqrt div def%
/deltaX0 tA DiffX0 mul rBs mul def
/deltaY0 tA rBs mul def
%-----------------------------------------------------------------------------------------------------------------------------------------
/Function ((x-X0+rBs*DiffX0/(sqrt(1+(DiffX0)^2))-rFs*(Diff)/(sqrt(1+(Diff)^2)))^2+%
            (func-funcX0+rFs/(sqrt(1+(Diff)^2))-rBs/(sqrt(1+(DiffX0)^2)))^2-dA12s^2) rpn def
%-----------------------------------------------------------------------------------------------------------------------------------------
/FunctionST ((x-XST+rBs*DiffxST/(sqrt(1+(DiffxST)^2))-rFs*(Diff)/(sqrt(1+(Diff)^2)))^2+%
            (func-funcxST+rFs/(sqrt(1+(Diff)^2))-rBs/(sqrt(1+(DiffxST)^2)))^2-dA12s^2) rpn def
%--------- inferior value to search for the intersection point of the frontwheel with the curve
/Zeros { %%  Funktion xinf
1 dict begin
/Xinf exch def % ---------------------------------- Untergrenze für die Schnittpunktsuche des Vorderrades mit der Kurve
% superior value to search for the intersection point of the frontwheel with the curve = X0 + distance axes + radius frontwheel + radius backwheel
/Xsup Xinf dA12s add rBs rFs add add def % ---------- Obergrenze für die Schnittpunktsuche des Vorderrades mit der Kurve
%---% Calculating the intersection point of the frontwheel with the function------ Schnittpunktberechnung ---------------------------------
/NB 0 def %-----------------loop-variable  -----  Laufvariable für loop
/NbreIterations 200 def % ---------- Maximum of iterations for the loop
{ %------------------------------------ loop begin ---------------------
 /xM Xinf Xsup add 2 div def %--------- Mittelwert von xM = (Xinf + Xsup):2
   /x Xinf def
   /F_1 FUNK def %----------------- F(Xinf)
   /x xM def %-------------------------
   /F_M FUNK def %----------------- F(xM)
   F_M 0 eq {exit} if %---------------- if F(xM) = 0 --> exit
    F_1 F_M mul 0 ge {/Xinf xM def} %-- F(Xinf) * F(xM) >= 0 Xinf = xM, else Xsup = xM
                     {/Xsup xM def}
    ifelse
 Xinf Xsup sub abs 1e-8 le {exit} if %- if Xinf - Xsup <= 10^-8 --> exit
  /NB NB 1 add def %------------------- else loopvariable NB = NB + 1  Loopvariable um eins erhöhen
 NB NbreIterations ge {exit} if %------ if number of iterations >= 200 --> exit
 } loop
xM
 end
} def
% \ifPst@EqValveStartPos
/FUNK {FunctionST} def XST Zeros /FWxST exch def
% \else /FWxST XST def \fi
/FUNK {Function} def X0 Zeros /FWx exch def
/FWy /x FWx def func def  %----------------------- Berührpunkt des Vorderrades (FWx,FWy)
/mFWy /x FWx def Diff def %----------------------- mFWy = Tangentensteigung in (FWx,FWy)
/TermFW 1 1 mFWy dup mul add sqrt div def %             1/sqrt(1+f'(x_Q)^2)
/deltaxFW TermFW mFWy mul rFs mul def %                  skal*rF*f'(x_Q)*1/sqrt(1+f'(x_Q)^2)
/deltayFW TermFW rFs mul def %                           skal*rF*1/sqrt(1+f'(x_Q)^2)
%------------------------------------------------------------------------------------
/KWRho {DiffI 1 Diff dup mul add 3 exp sqrt div} def
/dPhiB {1 rBs div KWRho sub dAB mul abs} def
/AngleCumB { %
 X1 X2 /x {dPhiB} eps SIMPSON
 } def
 /dPhiF {1 rFs div KWRho sub dAB mul abs} def
/AngleCumF { %
 X1 X2 /x {dPhiF} eps SIMPSON
 } def
% % length of the curve for the wheels  -- La longueur de la courbe pour la roue avant  --- Kurvenlänge
/X1 XST def /X2 X0 def %  Integral_{0}^{X0}
/sB AngleCumB def % ---backwheel - length from 0 to X0 ---- roue arrière ---- Kurvenlänge von 0 bis X0 (Hinterrad)
/X1 FWxST def /X2 FWx def
/sF AngleCumF def % --frontwheel - length from 0 to the abscissa of the intersection frontwheel with curve - roue avant -- Kurvenlänge von 0 bis zum SP Vorderrad - Kurve
%--------------------------------------------------------------------------------------------------------------------
%---% Definition angle of rotation for the backwheel --- Definition de l'angle de roue arriere ---- Rotationswinkel des Hinterrades
/phiB sB RadToDeg neg def
%---% Definition angle of rotation for the frontwheel --- Definition de l'angle de roue avant  ---- Rotationswinkel des Vorderrades
/phiF sF RadToDeg neg def
%---------------------------------------------------------------------------------------------------------------------
%--% coordinates for the axes of the wheels within the non-scaled system -- Koordinaten der Vorderradachse im nicht skalierten System ( außerhalb der \psscalebox )
%----------------------------------------------------------------------------------------------------------------------
/AFx FWx deltaxFW sub def %          % x-coordinate front axis   x-Koordinate der Vorderradachse
/AFy FWy deltayFW add def %          % y-coordinate front axis   y-Koordinate der Vorderradachse
/ABx X0 deltaX0 sub def %            % x-coordinate back axis    x-Koordinate der Hinterradachse
/ABy funcX0 deltaY0 add def %        % y-coordinate back axis    y-Koordinate der Hinterradachse
%--------------- Koordinaten der Vorderradachse im System des Fahrzeugs, also unskalierte Größen verwenden --( innerhalb der \psscalebox )
%--coordinates for the front axis of the wheels within \psscalebox  -- Die Koordinaten der Hinterradachse sind im System des Fahrzeugs  (0,0)
/AF1x dA12 dup mul rF rB sub dup mul sub sqrt def
/AF1y rF rB sub def
%------% slope and angle of the vehicle-----------------------------------------------------------------------------
/m-vehicle AFy ABy sub AFx ABx sub div def
/beta m-vehicle 1 atan def %          % angle for both radii of front- and backwheel are equal -- Neigungswinkel des Fahrzeugs bei gleich großen Rädern
/alpha AF1y neg AF1x atan def %       % additional angle if radii of front- and backwheel are not equal --  zusätzlicher Winkel bei unterschiedlich großen Rädern
% whole angle (correction with +180 in case the whole angle gets 90 degrees which can be possible with not equal radii)
/gamma beta alpha add AFx ABx lt { 180 add } if def %   gesamter Neigungswinkel des Fahrzeugs
%--------% Special case (mono-axis vehicle) Nr 4 --> segway ---------------------------------------------------------------------------
\ifPst@MonoAxis DiffX0 1 atan /omega exch def DiffX0 /mVehicle exch def \else FWy funcX0 sub FWx X0 sub div /mVehicle exch def mVehicle 1 atan /omega exch def \fi
/normNorm FWy funcX0 eq { 1 } { X0 FWx sub FWy funcX0 sub div dup mul 1 add sqrt } ifelse def
/mTgy FWy funcX0 eq { 1 } { X0 FWx sub FWy funcX0 sub div normNorm div } ifelse def
/mTgx FWy funcX0 eq { 0 } { 1 normNorm div } ifelse def
/xMTg X0 FWx add 2 div def
/yMTg funcX0 FWy add 2 div def
}%
%------------------------%%%%  END OF PS-CODE %%%%%--------------------------------------------------------
\pnode(!FWx FWy){FW}%
% % angle between the axes   angle de droite entre les axes
\rput{!gamma}(!ABx ABy){%  Das Fahrzeug wird mit Hilfe der Hinterradachse (ABx,ABy) und Gesamtdrehwinkel gamma gesetzt
%------------------------- % SETTING SOME VEHICLES  -------------------------------------------------------
\psscalebox{#1}{%---------------Das Fahrzeug kann skaliert werden -----------------------------------------
%\psset{linecolor=#2}%
\pnode(!\pst@GravNode){GravC}%
\pst@vehicle}}%
\ifPst@showSlope
\rput{!omega}(!X0 mVehicle 0 ge { 0 add } { 100 add } ifelse funcX0 0.75 sub){%
\psframebox[fillstyle=solid,fillcolor=greenSlp,linestyle=none]{\footnotesize\color{white}$m\geq 0$}}
\rput{!omega}(!X0 mVehicle 0 lt { 0 add } { 100 add } ifelse funcX0 0.75 sub){%
\psframebox[fillstyle=solid,fillcolor=redSlp,linestyle=none]{\footnotesize\color{white}$m < 0$}}
\ifPst@MonoAxis
\psplotTangent[linewidth=0.75\pslinewidth,linecolor=orange,nodesep=-2]{X0}{1}{#3}
\psplotTangent[linewidth=0.75\pslinewidth,linecolor=greenSlp,Tnormal,nodesep=-2]{X0}{1}{#3}
\psdot[linecolor=red](!X0 funcX0)
\else
\pcline[linecolor=greenSlp,nodesepA=-1,nodesepB=-3](!xMTg yMTg)(!xMTg mTgy 0 ge { mTgx } { mTgx neg } ifelse add yMTg mTgy abs add)
\pcline[linecolor=magenta,nodesep=-2](FW)(!X0 funcX0)
\psdot[linecolor=red](FW)
\psdot[linecolor=red](!X0 funcX0)
\fi%
\fi%
\showpointsfalse
\end@SpecialObj
\endgroup\ignorespaces}%
%
\def\SlopeoMeter#1#2{%
\colorlet{slpmColor}{#1}
\rput{0}(0,0){%
\pscircle[fillstyle=solid,fillcolor=black!90,linewidth=0.5pt,linestyle=solid](0,0){2.2}
\rput{!#2}(0,0){%
\multido{\r=.500+-.008,\rC=0+0.02}{25}{%
\pscircle[linewidth=0.008,linecolor=slpmColor,strokeopacity=\rC](0,0){!\r}
}
\multido{\i=0+2,\r=.400+-.008,\rC=0+0.02}{50}{%
\pswedge[linestyle=none,linewidth=0,fillstyle=solid,fillcolor=slpmColor,opacity=\r](0,0){2}{\i}{!\i\space 2 add}
}}
\multido{\nA=-90+10}{19}{%
\psline[linecolor=slpmColor](1.8;\nA)(1.95;\nA)
\rput(1.65;\nA){\psscalebox{0.4}{\color{slpmColor!10}\nA}}
}
\rput(-0.95,0.85){\psscalebox{0.6}{\color{slpmColor!20}\texttt{Steigungswinkel}}}
\rput(-0.95,0.5){\psscalebox{0.75}{\color{slpmColor!40}\texttt{Slope-o-Meter}}}
\rput{!#2}(0,0){%
\pspolygon[fillstyle=solid,fillcolor=slpmColor!50,linecolor=slpmColor,linejoin=2,linewidth=0.1pt](0.1;15)(0.1;-15)(1.75;0)
\pscircle[fillstyle=solid,fillcolor=black,linestyle=none,linewidth=0pt](0,0){0.3}
}}}
\catcode`\@=\PstAtCode\relax
\endinput
